\chapter{Conclusion and Future Work}
\label{chap:conclusion}

\section{Summary}
\label{sec:summary}

This thesis set out to measure how knowledge graphs change between releases, at
a scale where the graphs do not fit in memory. It answers its four questions as
follows.

\textbf{RQ1, can a metric catalogue be computed exactly through a SPARQL
endpoint?} Yes. Thirteen metrics across the class, entity and global level were
implemented in Rust and run against seven snapshots, the largest holding
2{,}489{,}858{,}800 triples. No metric loads the graph. Entropy is the case that
looks impossible at first, because the definition ranges over every object value
of a class, and Section~\ref{subsec:count_of_counts} shows it is not: entropy
depends only on the multiset of value counts, so a nested \texttt{GROUP BY}
folds hundreds of millions of values into a few hundred histogram rows and the
result is exact rather than sampled. The profile the catalogue produces behaves
as the construction of each graph predicts. Labels and \texttt{sameAs} carry the
most information for \texttt{Person}, gender carries almost none, and DBpedia's
predicate entropy exceeds YAGO's on every shared class because DBpedia draws
from 54{,}933 predicates where YAGO draws from 123.

\textbf{RQ2, how do the metrics change between versions, and where does the
change concentrate?} The change concentrates sharply, and the two graphs
concentrate it in opposite places. YAGO replaced 94\% of its class vocabulary
between 4 and 4.5.0.2, adding 111{,}122 classes while keeping 7{,}181, and held
its predicates nearly still at 72 kept. DBpedia did the reverse over a decade,
keeping 700 of 740 classes while 30{,}788 predicates disappeared and 23{,}897
arrived. At class level, metric~9 localises YAGO's change to the content that
left: \texttt{GeoCoordinates} reached churn 1.0 as pure deletion, and
\texttt{Article} recorded 758{,}368 deletions against 17 additions, which
accounts for most of the 1.18 billion triple drop. \texttt{Person} behaved
differently, 223{,}227 added against 245{,}398 deleted, which is revision rather
than retirement. Separating those two cases is what the added and deleted split
provides and a single churn number cannot.

\textbf{RQ3, can version comparison run on commodity hardware?} Yes, and the
integer encoding is what makes it so. Against a string-based baseline on
identical inputs, the encoded path ran 2.34 to 6.17 times faster and held its
working set in 16 to 19 times less memory. The memory result is the one that
decides feasibility: extrapolating the baseline's 575\,MB for three million
triples to a hundred-million-triple comparison gives roughly 19\,GB, which does
not fit on the 16\,GB client, while the encoded path's 35\,MB extrapolates to
about 1.2\,GB, which does. Both paths run inside the same invocation and the
implementation asserts they return identical sets, so the speedup cannot drift
away from being correct.

\textbf{RQ4, what do the metrics say across two different graphs?} They work, and
they report a consistent structural difference. Metric~13 matched sixteen
classes automatically between YAGO~4.5.0.2 and DBpedia 2025 and found DBpedia
using at least an order of magnitude more predicates on every one, with ratios
from 24 to 92: 3{,}025 against 31 on \texttt{building}, 2{,}888 against 34 on
\texttt{administrativearea}, 1{,}315 against 24 on \texttt{mountain}. Sixteen
classes without an exception makes that structural.

The exception is worth more than the rule. \texttt{Organization} is the only
class of the sixteen where the two graphs report the same entropy, 14.2752 bits
against 14.2754. Repeating the comparison against DBpedia 2015, 2016 and 2022
shows the agreement is recent: across those three releases the two graphs sat
between 0.63 and 0.72 bits apart, and the distance vanished only between 2022
and 2025, while YAGO's value never moved. Over the same decade the graphs ended
further apart on nine of the fifteen classes measured in all four releases, so
this is an exception against a trend. Repeating the series with YAGO~4 on the
YAGO side shows the agreement needed both graphs to move: against YAGO~4,
\texttt{organization} stays 4.5 to 5.9 bits from DBpedia in every release,
and YAGO's own value fell by 5.19 bits between its two releases. It is the
single result in this thesis I did not expect, and four time points are not
enough to explain it.

The entropy-weighted PageRank variant of
Section~\ref{sec:entropy_weighted_pagerank} moved 5{,}327 of 5{,}372 sampled
nodes, which makes it a different ranking rather than a reordering of ties. On
the sampled YAGO~4 subgraph it converged in 43 iterations where the unweighted
baseline did not converge in 100.

\section{Limitations}
\label{sec:limitations}

Five limitations constrain these results, and each is a consequence of a design
decision made deliberately rather than an oversight.

\textbf{The triple diff is exact inside a subject window, not over the whole
graph.} Metric~10 takes the lexicographically first $n$ subjects from each side
and cuts both at the smaller bound. Inside that window the result is exact. It
says nothing about subjects outside it. The alternative, sorting both graphs
globally, is what the subject window exists to avoid: QLever asked for 12.7\,GB
to sort a graph far smaller than YAGO and refused the query. Whole-graph windows
on billion-triple graphs are not feasible through SPARQL, and the implementation
says so rather than returning a number that looks complete.

\textbf{The YAGO~4 cross-graph runs leave out two root classes.} YAGO~4's class
distribution is far more top-heavy than YAGO~4.5's: its largest classes are
\texttt{Thing} at 66.9 million entities and \texttt{CreativeWork} at 35.9
million, where YAGO~4.5's largest is \texttt{Taxon} at 3.5 million. Both match
a DBpedia class by name. Counting their predicates became feasible once the
query grouped by predicate instead of asking for \texttt{COUNT(DISTINCT)}, which
had requested 11.1\,GB, but the entropy histogram still pushed the QLever server
past the 12\,GB of the machine hosting it, and the operating system killed it.
Both classes are therefore excluded from every YAGO~4 run, which the result files
record. The loss is small, because a root class measures the whole graph rather
than a concept, but it is a gap that more memory would close. The
\texttt{organization} result also still rests on one class and four time points,
now measured against two YAGO releases.

\textbf{Cross-graph class matching is by local name.} Metric~13 treats
\texttt{schema:Person} and \texttt{dbo:Person} as the same class because the
local names agree. It does not check that they contain the same entities, and where a graph
carries several aliases of one class the choice among them is arbitrary: DBpedia
2025 holds four \texttt{Person} classes within six entities of one another, and
the metric matched one of them without preferring it. The method also yields only
fifteen or sixteen matched classes from 150 candidates per side for YAGO~4.5,
and twenty-three to twenty-five for YAGO~4, because two differently built
vocabularies rarely agree on a local name. Ringler
and Paulheim~\cite{ringler2017one} show these graphs are not interchangeable, so
the assumption is known to be imperfect, and a $\Delta H$ of 6.49 bits on
\texttt{person} compares two class extents that share a name rather than two
descriptions of the same entities.

\textbf{Per-class metrics cover the top classes, not all of them.} A snapshot
with 118{,}303 classes cannot have every class measured at reasonable cost, so
the runs cover the most populated ones. The pinned pass mitigates this for
comparison by fixing the class list, but a change confined to rare classes would
not appear in these results.

\textbf{YAGO~3 cannot be compared per class to YAGO~4 at all.} The two releases
share no class names, because YAGO~3 types entities against WordNet synsets and
Wikipedia categories under a four-layer taxonomy while YAGO~4 uses schema.org. A
per-class comparison would need a hand-built synset mapping and a transitive
\texttt{rdf:type/\allowbreak rdfs:subClassOf*} closure, and without them metric~10 over that
pair reports 56{,}021 additions and zero deletions: correct arithmetic, no
meaning. The aggregate metrics still capture the overhaul, with the class count
falling from 475{,}673 to 10{,}103.

\textbf{Entropy-weighted PageRank normalises by out-degree, not by weight mass.}
Each edge contributes $w(p) \cdot \mathrm{PR}(u) / \mathrm{outdeg}(u)$, so a node
whose edges all carry low weights leaks score rather than redistributing it. The
ranking is still a valid ordering, but it is not a probability distribution in
the way classic PageRank is. The results of Section~\ref{sec:exp4} also come from
a sampled subgraph of 5{,}372 nodes, not the full graph, so the convergence
observation is a single data point and should not be read as a general claim.

\section{Future Work}
\label{sec:future_work}

Four directions follow.

\textbf{More graphs, starting with Wikidata.} Wikidata~\cite{vrandecic2014wikidata}
is the obvious third graph, and since YAGO~4 already uses Wikidata identifiers
the class matching of metric~13 would have something better than name equality
to work with. That alone would address the largest limitation above.

\textbf{Instance-level alignment for cross-graph comparison.} Replacing name
matching with \texttt{owl:sameAs} linkage would turn metric~13 from a comparison
of similarly named classes into a comparison of the same entities, and would let
the \texttt{organization} result be tested rather than only observed. The
links would need checking first: Halpin et al.~\cite{halpin2010sameas} show that
\texttt{owl:sameAs} is often used for similarity rather than strict identity,
which would reintroduce the imprecision the alignment is meant to remove.

\textbf{Pushing the diff beyond the window.} The integer encoding is not what
limits metric~10; the need to avoid a global sort is. A sorted index over
subjects, or a partitioned comparison that sweeps the subject space in windows
and accumulates, would extend exactness to the whole graph at a cost that the
memory figures in Section~\ref{sec:exp3} suggest is affordable.

\textbf{Streaming metric updates.} Every metric here recomputes from scratch. For
graphs published as incremental changesets, updating a metric from the changeset
rather than the full graph would make continuous monitoring practical, which is
what a consumer tracking a graph across releases actually wants.

Finally, the YAGO~3 synset mapping would close the one comparison this thesis
cannot make. It is a bounded piece of work, a mapping table plus a transitive
closure at index time, and it would extend the YAGO evolution series back one
full generation to cover the ontology replacement at class level rather than only
in aggregate.
